\documentclass[english,a4paper,12pt]{article}
%\documentclass[english,journal=jctcce,manuscript=article,articletitle=true,layout=twocolumn]{achemso}
\usepackage[english]{babel}
\usepackage[T1]{fontenc}
\usepackage[margin=2cm]{geometry}

\usepackage{doi}
\usepackage[version=4]{mhchem}
\usepackage{multicol}

\title{}
\author{}
\date{}

\begin{document}

\pagestyle{empty}

\section*{Theory}

\noindent The ADAPT ansatze is
%
\begin{equation}
|\psi(\mathbf{\theta})\rangle_N=\prod_{k=1}^{N}e^{\theta_k}\hat{\tau}_k|\psi_\mathrm{ref}\rangle.
\end{equation}
%
At each iteration we grow the ansatze as
%
\begin{equation}
|\psi(\mathbf{\theta})\rangle_{N+1}=e^{\theta_{N+1}\hat{\tau}_{N+1}}|\psi(\mathbf{\theta})\rangle_N,
\end{equation}
%
where
%
\begin{equation}
\begin{split}
\hat{\tau}_{N+1}&=\text{argmax}_{\hat{\tau}_i\in\{A_j\}_n}\ \frac{\partial}{\partial\theta_i}\langle E\rangle \\
&=\text{argmax}_{\hat{\tau}_i\in\{A_j\}_n}\ _N\langle\psi(\mathbf{\theta})|[\hat{H},\hat{\tau}_i]|\psi(\mathbf{\theta})\rangle_N.
\end{split}
\end{equation}
%
$\{A_j\}_n$ is the operator pool for $n$ qubits.
We have two options:
%
\begin{enumerate}
\item $\{A_j\}_n=\{V_j\}_n=\{Z_n\{V_k\}_{n-1}, iY, iY_{n-1}\}$ and $\{V_j\}_2=\{iZ_2Y_1,iY_2\}$.
\item $\{A_j\}_n=\{G_j\}_n=\{iZ_{k+1}Y_k,iY_{k+1}\}_{k=0}^{n-1}$.
\end{enumerate}
%
The second results in only local two-qubit operations.
The gradients are measured on the quantum device.
At the end of each iteration we optimize $\mathbf{\theta}$ classically.
Convergence of the algorithm is achieved when
%
\begin{equation}
\vert\vert g\vert\vert=\left(\sum_i\left(\frac{\partial}{\partial\theta_i}\langle E\rangle\right)^2\right)^{1/2}\leq\epsilon.
\end{equation}
%

\section*{Practical stuff}

The reference state $|\psi_\mathrm{ref}\rangle$ is computed with the \textsc{Lucas} code.
We can begin with the Li and Be atoms and if they work we can move to heavier atoms and/or simple molecules such as LiH and H$_2$O.

\paragraph{Questions}

\begin{enumerate}
\item Perform \textsc{Lucas} calculation
\item Read out the reference state
\item Choose the operator pool
\item Implement the state on Q50
\begin{itemize}
\item How is $|\psi_\mathrm{ref}\rangle$ mapped to Q50?
\end{itemize}
\item Measure the gradients on Q50
\begin{itemize}
\item How are the measurements performed?
\end{itemize}
\item Select largest gradient
\item Readout
\begin{itemize}
\item How do we read the gradients from Q50?
\end{itemize}
\item Reoptimize classically
\end{enumerate}

\end{document}